\section{Metodología}\label{sec:metodologia}

\subsection{Enfoque metodológico}

El diseño metodológico del presente trabajo se basa en la construcción de un benchmarking sistemático para determinar cómo las decisiones en la etapa de preparación de datos condicionan el rendimiento de modelos de aprendizaje supervisado en el dominio clínico.

\subsection{Evaluación y selección de datasets}

\textbf{Objetivo de la evaluación inicial}

La fase inicial de selección de datos constituyó una etapa crítica en la
configuración metodológica. El objetivo primordial no fue simplemente adquirir un conjunto de datos biomédicos, sino identificar un entorno de información que presentara una complejidad intrínseca y desafíos en el preprocesamiento, adecuados para validar la hipótesis central
del estudio (Sección~\ref{sec:hipotesis}). Para ello, se requería un conjunto de datos que reflejara la heterogeneidad y los desafíos de los datos reales en salud, incluyendo valores faltantes, informativos, multicolinealidad fisiológica y el potencial de fuga de datos
temporal o estructural.

\textbf{Datasets evaluados y desestimados}

Breast Cancer Wisconsin (Diagnostic)~\citep{wolberg1993}

\begin{itemize}
  \item Descripción breve: Este dataset, ampliamente utilizado en la comunidad de
    Machine Learning, se centra en la clasificación de tumores mamarios como
    benignos o malignos. Contiene 569 observaciones y 32 variables, la mayoría
    de las cuales son medidas biomédicas numéricas continuas (e.g., radius\_mean,
    texture\_mean, perimeter\_mean) derivadas de imágenes digitalizadas de masas
    celulares. La variable objetivo es el diagnóstico clínico (M para maligno, B
    para benigno).
  \item Ventajas observadas: Es un dataset excepcionalmente "limpio", con una
    documentación exhaustiva y es un referente estándar para el benchmarking de
    algoritmos de clasificación. La ausencia de valores nulos significativos
    simplifica su uso.
  \item Limitaciones observadas: A pesar de sus ventajas, su baja complejidad
    clínica y estructural lo hizo inadecuado para los objetivos de este
    estudio. La escasez de valores faltantes (no se observaron valores nulos en
    ninguna de las 569 observaciones) limita drásticamente la oportunidad de
    explorar estrategias avanzadas de imputación basadas en conocimiento
    clínico. Además, al consistir principalmente en variables numéricas ya
    procesadas, ofrece un espacio reducido para el análisis profundo de la
    heterogeneidad de variables categóricas, económicas o de gestión
    hospitalaria, desviándose de la intención de demostrar la robustez del
    preprocesamiento en escenarios más realistas.
\end{itemize}

Parkinsons~\citep{little2007}

\begin{itemize}
  \item Descripción breve: Se evaluó el dataset parkinsons.data, el cual contiene
    195 observaciones y 24 variables biomédicas derivadas de señales de voz,
    destinadas a clasificar la presencia de la enfermedad de Parkinson. La
    variable objetivo es status (0 para sanos, 1 para pacientes).
  \item Ventajas observadas: Representa un problema clínico relevante y utiliza
    señales biomédicas reales.
  \item Limitaciones observadas: Su principal limitación para este estudio fue su
    tamaño extremadamente reducido. Con tan solo 195 registros, la capacidad de
    generalización y la robustez de los modelos entrenados se verían
    comprometidas, especialmente en un contexto de benchmarking con múltiples
    estrategias. Además, presenta una menor heterogeneidad de variables en
    comparación con un entorno hospitalario complejo, lo que reduce el alcance
    para explorar estrategias avanzadas de imputación o selección de variables
    que trasciendan lo puramente estadístico. La relativa "perfección" de sus
    datos y la menor riqueza clínica y hospitalaria también lo desvían del
    propósito de este TFM.
\end{itemize}

Chronic Kidney Disease (CKD)~\citep{rubini2015}

\begin{itemize}
  \item Descripción breve: Este dataset busca clasificar la presencia de enfermedad
    renal crónica (ckd o notckd) en pacientes, a partir de 400 observaciones
    y 25 variables clínicas y biomédicas, que incluyen parámetros fisiológicos,
    análisis sanguíneos y condiciones clínicas del paciente.
  \item Ventajas observadas: Representa un problema clínico relevante y exhibe una
    mezcla de variables numéricas y categóricas, junto con una cantidad
    considerable de valores faltantes en múltiples variables (ej., rbc, rbcc,
    wbcc, pot).
  \item Limitaciones observadas: A pesar de ser más prometedor en términos de
    presencia de NAs y diversidad de variables, su tamaño limitado (400
    observaciones) y su menor diversidad clínica en comparación con escenarios
    hospitalarios de cuidados críticos lo hicieron menos adecuado. La escala del
    dataset no permitía una exploración tan exhaustiva de la interacción entre
    variables complejas o la validación de estrategias de preprocesamiento
    avanzadas con la solidez requerida para un estudio de benchmarking integral.
\end{itemize}

\textbf{Criterios de evaluación y razones de la desestimación}

Los datasets evaluados, si bien válidos para propósitos académicos específicos,
fueron finalmente descartados por presentar limitaciones críticas para los
objetivos de este trabajo. Los criterios de evaluación considerados fueron:

\begin{itemize}
  \item Complejidad Clínica y Estructural: La capacidad del dataset para reflejar la
    interconexión de múltiples factores clínicos, demográficos y
    administrativos.
  \item Diversidad de Variables: La presencia de una mezcla amplia de variables
    continuas, categóricas, demográficas y administrativas.
  \item Tamaño del Dataset: Un número suficiente de registros para asegurar la
    robustez estadística de los modelos y la estabilidad de las particiones.
\end{itemize}

Estos factores reducirían la capacidad de demostrar el
impacto diferencial de un preprocesamiento avanzado y la comprensión de las
decisiones tomadas bajo un razonamiento clínico-estadístico.

\subsection{Selección del dataset SUPPORT2}

\textbf{Justificación de selección y variable objetivo}

El dataset SUPPORT2 (Study to Understand Prognoses and Preferences for Outcomes
and Risks of Treatments)~\citep{knaus_support2} fue seleccionado como el eje central de este estudio
por su alta dimensionalidad, riqueza de información y su fiel representación
de los desafíos inherentes a los datos hospitalarios del mundo real. Este
dataset no solo ofrece una alta dimensionalidad con 9,105 registros y 47
variables, sino que también presenta patrones de datos faltantes significativos
y una profunda interconexión de variables que demandan un análisis exploratorio
y un preprocesamiento rigurosos.

\textbf{Justificación de la Variable Objetivo: hospdead}

La variable hospdead (mortalidad hospitalaria) fue elegida como la variable
objetivo principal para el problema de clasificación binaria. Esta decisión se
fundamenta en un análisis detallado que considera la naturaleza del
dataset, su contexto clínico original y las observaciones clave del EDA:

\begin{enumerate}
  \item Coherencia con el Diseño del Dataset: La documentación original de SUPPORT2
    y su propósito principal se centran en la predicción del pronóstico y los
    outcomes en pacientes críticos. hospdead representa un desenlace directo y
    central a esta misión, caracterizando un evento bien definido dentro del
    período de hospitalización.
  \item Relevancia Clínica y Metodológica: La mortalidad intrahospitalaria es una
    métrica de alto interés biomédico, fundamental para la predicción de riesgo, la toma de decisiones clínicas y la evaluación de
    la severidad del paciente. Al integrar múltiples variables
    fisiológicas, funcionales y clínicas agudas, está especialmente diseñado
    para modelar la evolución del paciente durante la hospitalización, haciendo
    de hospdead un outcome natural y congruente.
  \item Estabilidad frente a Alternativas:
    \begin{itemize}
      \item death (mortalidad general): Aunque también representa mortalidad, death
        abarca un horizonte temporal más amplio (hasta seis meses post-ingreso,
        incluyendo eventos fuera del hospital) y puede incorporar mayor pérdida
        de seguimiento y factores externos. hospdead ofrece una variable
        objetivo más controlada y estable para un problema de clasificación
        supervisada.
      \item Variables Diagnósticas (dzgroup, dzclass): Transformar estas variables
        categóricas complejas en un problema binario (ej., "cáncer vs. no
        cáncer") implicaría agrupar condiciones clínicas dispares bajo una
        clase negativa, lo que reduciría la coherencia biológica y la
        discriminación inherente al dataset.
      \item Condiciones Crónicas (diabetes, dementia): Estas variables, si bien
        importantes, se relacionan más con características preexistentes del
        paciente que con la evolución clínica aguda en el hospital, que es el
        foco principal de SUPPORT2 y de muchas de sus variables fisiológicas.
    \end{itemize}
\end{enumerate}

Durante el EDA, se observaron patrones de relación intrigantes entre hospdead y
una amplia gama de predictores (histogramas, boxplots, correlaciones). Estos
patrones sugirieron que hospdead es una variable rica en información, capaz de
revelar cómo las interacciones complejas de factores clínicos y administrativos
influyen en el pronóstico, lo que la convierte en una elección óptima para
construir un modelo predictivo robusto.

\textbf{Descripción general del dataset}

El dataset SUPPORT2 es el resultado de un estudio multicéntrico observacional
que recopiló información clínica detallada de 9,105 pacientes adultos ingresados
en cinco hospitales de Estados Unidos entre 1989 y 1994. Estos pacientes
padecían enfermedades graves en fase avanzada, incluyendo insuficiencia
respiratoria aguda, EPOC, insuficiencia cardíaca congestiva, cirrosis, cáncer,
coma y fallo multiorgánico. El objetivo principal del estudio fue analizar el
pronóstico, la evolución clínica y los outcomes hospitalarios, especialmente la
mortalidad y las preferencias de atención al final de la vida. La magnitud y
profundidad de los datos recopilados lo establecen como un recurso de referencia
en la investigación de outcomes en pacientes críticos.

\textbf{Tipos de variables presentes}

La riqueza del dataset radica en su diversidad de 47 variables, que abarcan
múltiples dominios. El diccionario completo de variables se presenta en el
Anexo~\ref{sec:anexo-diccionario}.

\begin{itemize}
  \item Variables Demográficas: Edad (age), sexo (sex), etnia (race), nivel de
    ingresos (income) y nivel educativo (edu).
  \item Variables Fisiológicas y Analíticas: Constantes vitales como la presión
    arterial media (meanbp), frecuencia cardíaca (hrt), frecuencia respiratoria
    (resp) y temperatura (temp). También incluye marcadores de laboratorio como
    la albúmina sérica (alb), bilirrubina (bili), creatinina (crea), glucosa
    (glucose) y pH sanguíneo (ph).
  \item Scores Clínicos Validados: Índices de severidad como el APACHE III (aps), el
    TISS (avtisst), la puntuación de coma basada en la escala de Glasgow
    (scoma), y la puntuación de fisiología SUPPORT (sps).
  \item Variables de Gestión y Económicas: Días de estancia hospitalaria (slos),
    cargos facturados por el hospital (charges) y costos totales calculados
    (totcst, totmcst).
  \item Variables Funcionales: Índices de Actividades de la Vida Diaria (ADL)
    reportados por el paciente (adlp) y por un familiar (adls).
  \item Variables Derivadas y de Pronóstico: Estimaciones de probabilidad de
    supervivencia a 2 (surv2m, prg2m) y 6 meses (surv6m, prg6m).
\end{itemize}

\subsection{Análisis Exploratorio de Datos (EDA)}

El análisis exploratorio reveló una imagen detallada de la estructura y
características del dataset, identificando patrones cruciales para las
etapas posteriores.

\textbf{Análisis de valores faltantes (Missing Values)}

La integridad de los datos en SUPPORT2 es un aspecto central, con una
distribución de valores nulos que varía desde el 0\% hasta un
considerable 61.9\%. Este patrón de missingness no es aleatorio, sino que es
altamente informativo:

\begin{itemize}
  \item Variables con Alta Ausencia de Datos (Missingness > 40\%):
    \begin{itemize}
      \item adlp (Actividades de la vida diaria previas): 61.95\% de valores
        faltantes (5,641 registros). Crucialmente, los pacientes fallecidos en
        el hospital presentaron un 92\% de nulos en esta variable, frente al 52\%
        de los supervivientes. Esto subraya que la ausencia del dato está
        directamente vinculada con la gravedad extrema del paciente (coma,
        sedación, imposibilidad de responder), y es una señal clínica relevante.
      \item urine (Producción urinaria): 53.40\% de valores faltantes (4,862
        registros).
      \item glucose (Glucosa sérica): 49.42\% de valores faltantes (4,500 registros).
      \item bun (Nitrógeno ureico en sangre): 47.80\% de valores faltantes (4,352
        registros).
    \end{itemize}
  \item Variables con Moderada Ausencia de Datos (Missingness 5\% - 40\%):
    \begin{itemize}
      \item alb (Albúmina sérica): 37.03\% (3,372 registros).
      \item bili (Bilirrubina sérica): 28.57\% (2,601 registros).
      \item pafi (Índice de oxigenación Kirby): 25.54\% (2,325 registros). Se observó
        una menor ausencia en pacientes con patologías respiratorias (ej.,
        un 13.7\% en COPD), sugiriendo una recolección condicionada al protocolo
        clínico.
      \item ph (pH sanguíneo): 25.09\% (2,284 registros). Similar a pafi, con menor
        ausencia en pacientes críticos/respiratorios.
      \item income (Nivel de ingresos): 32.75\%.
      \item adls (Actividades de la vida diaria al seguimiento): 31.49\%. A
        diferencia de adlp, sus patrones de missingness son menos influenciados
        por la gravedad extrema.
    \end{itemize}
  \item Variables con Mínima Ausencia de Datos (Missingness < 5\%):
    \begin{itemize}
      \item Variables vitales como sod (0.01\%), temp (0.01\%), resp (0.01\%), hrt
        (0.01\%), meanbp (0.01\%), surv6m (0.01\%), surv2m (0.01\%), aps (0.01\%),
        sps (0.01\%), scoma (0.01\%). Estas variables, con datos casi completos,
        constituyen el bloque más fiable del dataset.
    \end{itemize}
\end{itemize}

\textbf{Estadísticas descriptivas y distribuciones de variables}

El análisis de las distribuciones y estadísticas descriptivas reveló una
heterogeneidad considerable:

\begin{itemize}
  \item Sesgo Positivo: Una gran mayoría de variables continuas, incluyendo
    crea, bun, bili, wblc, charges, totcst, totmcst, slos (estancia
    hospitalaria) y hday (día de ingreso), muestran una fuerte asimetría
    positiva. Esto se manifiesta con una acumulación de registros en los valores
    bajos y una cola muy larga hacia la derecha. Por ejemplo, slos tiene una
    mediana de 11 días pero un máximo de 343 días, y charges oscila desde una
    mediana de 25,024 hasta un máximo de 1,435,423. Esto refleja una
    población con un estado relativamente estable, pero con un subgrupo de
    casos extremos que disparan los costos y la duración de la estancia.
  \item Comportamientos Bimodales: Variables como la frecuencia cardíaca (hrt),
    presión arterial media (meanbp) y temperatura (temp) exhiben dos picos
    distintos. Esto indica la presencia de dos subpoblaciones claras: pacientes
    con valores dentro de rangos fisiológicos normales y pacientes
    descompensados (ej., taquicardia/bradicardia, hipotensión/hipertensión,
    fiebre/hipotermia).
  \item Comportamientos Multimodales: Las variables de estimación subjetiva del
    pronóstico (prg2m, prg6m) muestran picos espaciados uniformemente
    (25\%, 50\%, 75\%). Este patrón sugiere un sesgo de redondeo por parte de los
    médicos al asignar probabilidades, lo que difiere de una variable biológica
    continua.
  \item Distribución Simétrica/Normal: Variables como la edad (age), el score de
    fisiología SUPPORT (sps), el pH sanguíneo (ph) y la concentración de sodio
    (sod) tienden a una distribución más simétrica o aproximadamente normal.
  \item Variable Objetivo (hospdead): Se observó un desbalance moderado, con 6,745
    pacientes (74.1\%) que sobrevivieron y 2,360 (25.9\%) que fallecieron durante
    la hospitalización.
\end{itemize}

\textbf{Análisis de correlación y detección de fuga de datos}

La matriz de correlación reveló importantes relaciones lineales y potenciales
fuentes de fuga de datos:

\begin{itemize}
  \item Redundancia Alta: Se identificó una correlación perfecta (r = 1.0) entre
    adlsc y adls, indicando que representan la misma información o
    transformaciones idénticas. Similarmente, charges, totcst y totmcst muestran
    correlaciones muy elevadas (superiores a 0.86), apuntando a una fuerte
    redundancia en la información económica.
  \item Leakage Temporal: El tiempo total de seguimiento del paciente (d.time)
    exhibe una marcada correlación inversa de -0.79 con death. Esta fuerte
    relación se debe a que el d.time se trunca abruptamente cuando ocurre el
    fallecimiento, lo que lo convierte en una variable que incorpora información
    del futuro respecto a una predicción al ingreso.
  \item Relación de Variables Derivadas: Las probabilidades estimadas de
    supervivencia a corto y mediano plazo (surv2m, surv6m) se encuentran
    fuertemente vinculadas de manera negativa con el score de severidad
    fisiológica (sps), alcanzando correlaciones de hasta -0.72. Esto sugiere
    que estas variables son, en parte, el resultado de modelos predictivos y no
    mediciones directas.
  \item Consistencia Fisiológica: Biomarcadores de función renal como el nitrógeno
    ureico (bun) y la creatinina (crea) exhiben una correlación positiva
    moderada-alta de 0.66, confirmando una progresión conjunta ante el deterioro
    del sistema urinario.
\end{itemize}

\subsection{Construcción de estrategias de preprocesamiento}

Esta sección detalla las dos aproximaciones metodológicas diseñadas para el
preprocesamiento del dataset SUPPORT2. La comparación de estas estrategias
constituye el núcleo del experimento, permitiendo evaluar si la integración del
conocimiento de dominio mejora la capacidad predictiva y la robustez de los
modelos de aprendizaje automático frente a un enfoque puramente algorítmico.

\subsubsection{Estrategia clínico-estadística (V1)}

\textbf{Filosofía y Razonamiento de Dominio}

La estrategia V1 se diseñó bajo un enfoque de supervisión experta en el proceso
de preprocesamiento. El razonamiento central es que, en un entorno hospitalario,
la ausencia de una medición o la correlación entre variables puede responder a
una lógica clínica, administrativa o temporal. Por lo tanto, la preparación de
los datos debe respetar la coherencia biológica. En lugar de aplicar reglas de
corte genéricas, como eliminar variables por un porcentaje fijo de nulos, se
analizó el origen de cada variable y el motivo de su ausencia para decidir su
permanencia o su forma de imputación.

\textbf{Gestión de la Integridad y Calidad del Dato}

Una premisa clave de esta estrategia es priorizar un conjunto de predictores
clínicamente consistente frente a la retención indiscriminada de variables con
alta incertidumbre o riesgo metodológico.

\begin{itemize}
  \item Eliminación de registros por inconsistencia basal: Se decidió eliminar las
    observaciones incompletas en variables como crea, meanbp, hrt, resp, temp,
    sod, scoma, aps y crea. Al presentar un porcentaje de nulos mínimo (<1.2\%),
    su eliminación supone una pérdida insignificante de registros pero garantiza
    que el modelo solo aprenda de pacientes con un perfil clínico basal
    completo. Además, se comprobó que estos nulos se concentraban en la clase
    mayoritaria, por lo que no se afectó el desbalance del dataset.
  \item Exclusión de variables por baja fiabilidad: Variables como bun y glucose
    fueron eliminadas debido a que sus altísimos niveles de nulos (cercanos
    al 50\%) obligarían a realizar imputaciones masivas que homogeneizarían
    artificialmente la muestra, destruyendo el sentido biológico del dato.
\end{itemize}

\textbf{Prevención de fuga de datos}

\begin{itemize}
  \item Leakage de Modelos Previos: Las variables surv2m y surv6m fueron eliminadas
    por ser estimaciones generadas por otros modelos predictivos, lo que
    inflaría la precisión de forma artificial. Asimismo, basándose en el
    protocolo original de SUPPORT (JAMA, 1995), se identificó que prg2m y prg6m
    formaron parte de una intervención clínica posterior, lo que las inhabilita
    como predictores de línea base.
  \item Fuga temporal: Se eliminó d.time (tiempo de seguimiento) y slos (estancia hospitalaria) por su
    separación matemática perfecta con el desenlace hospdead, y se descartaron
    sps y avtisst por representar métricas procesadas que pueden verse alteradas
    por decisiones terapéuticas posteriores al ingreso.
\end{itemize}

\textbf{Imputación Basada en Contexto y Estratificación}

En lugar de una mediana global, se aplicó una imputación estratificada por
diagnóstico (dzgroup), bajo la lógica de que los valores normales varían según
la patología:

\begin{itemize}
  \item Variables Fisiológicas Condicionadas: Para bili, alb, ph, pafi y wblc, se
    utilizó la mediana del grupo diagnóstico. Es fisiológicamente coherente que
    un paciente con cirrosis tenga una mediana de bilirrubina distinta a uno con
    insuficiencia cardíaca.
  \item Análisis del Blanco Analítico: En el caso de wblc (glóbulos blancos), se
    validó que su ausencia no dependía de estancias cortas (la mediana de
    estancia para estos nulos fue de 9 días), lo que justifica una imputación
    por contexto en lugar de una eliminación.
  \item Incertidumbre Administrativa: Para race, dnr, dnrday y edu, se imputó la
    categoría "Unknown". Se asume que la falta de registro en estas variables
    puede reflejar aspectos administrativos o del proceso de recolección que
    tienen valor predictivo latente.
\end{itemize}

\textbf{Interpretación de la Ausencia Informativa (ADL)}

Uno de los puntos más críticos fue el tratamiento de las escalas funcionales:

\begin{itemize}
  \item Información Oculta en Nulos: En la variable adls, se sustituyeron los nulos
    por la categoría 8 ("No completado"). Se interpreta que el hecho de no poder
    completar el formulario está relacionado con la gravedad o el estado de
    conciencia del paciente.
  \item Exclusión de adlp por Sesgo de Recolección: Se observó que el 92\% de los
    fallecidos no tenían datos en adlp frente al 52\% de los vivos. Aunque tenía
    un alto poder predictivo, se determinó que este poder era un "artefacto" del
    proceso de recolección (los pacientes más graves no podían responder) y no
    una característica fisiológica. Para priorizar la interpretabilidad clínica,
    se decidió eliminarla.
\end{itemize}

\textbf{Refinamiento Estructural}

\begin{itemize}
  \item Eliminación por Redundancia y Control Metodológico: Se eliminó dzclass por
    ser una agrupación redundante de dzgroup, y adlsc por ser una versión ya
    imputada previamente en el estudio original sobre la cual se perdió el
    control metodológico.
  \item Simplificación Económica: Se eliminaron charges (con solo un 2\% de
    nulos), totcst, totmcst e income por su alta correlación y
    baja relevancia clínica directa frente a la mortalidad aguda, además de que
    son variables que no están disponibles al ingreso y pueden introducir fuga
    de datos temporal.
\end{itemize}

Este conjunto de decisiones asegura que V1 no solo sea
estadísticamente apto, sino clínicamente veraz, sentando las bases para un
benchmark donde el conocimiento de dominio sea el factor diferencial.

\begingroup
\footnotesize
\setlength{\tabcolsep}{3pt}
\begin{longtable}{@{}p{0.23\textwidth}p{0.17\textwidth}p{0.21\textwidth}p{0.31\textwidth}@{}}
\caption[Preprocesamiento V1]{Preprocesamiento clínico-estadístico (V1)}\label{tab:preproc-v1}\\
\toprule
\textbf{Variable} & \textbf{Acción} & \textbf{Estrategia aplicada} & \textbf{Observación / Justificación} \\
\midrule
\endfirsthead
\multicolumn{4}{c}{\tablename\ \thetable\ — V1 (continuación)} \\
\toprule
\textbf{Variable} & \textbf{Acción} & \textbf{Estrategia aplicada} & \textbf{Observación / Justificación} \\
\midrule
\endhead
\bottomrule
\endfoot
meanbp, hrt, resp, temp, sod, scoma, aps, crea & Eliminar observaciones & Eliminación puntual & Nulos mínimos en variables basales; se priorizó conservar registros clínicamente completos. \\
race, dnr, edu & Imputar & Categoría ``Unknown'' & La ausencia se trató como información administrativa potencialmente útil. \\
wblc, bili, alb, ph, pafi & Imputar & Mediana por dzgroup & Imputación condicionada por diagnóstico para preservar coherencia fisiológica. \\
adls & Imputar & Valor 8 & Se codificó como ``No completado'', interpretando la ausencia como señal clínica. \\
urine, bun, glucose & Eliminar variable & Missing elevado & La imputación masiva podía homogeneizar artificialmente variables clínicas relevantes. \\
prg2m, prg6m, surv2m, surv6m & Eliminar variable & Riesgo de leakage & Variables pronósticas o derivadas de modelos previos, no adecuadas como predictores basales. \\
death, d.time & Eliminar variable & Leakage directo y temporal & Incorporan información del desenlace o del seguimiento posterior. \\
slos, charges & Eliminar variable & Información posterior & Estancia hospitalaria y cargos acumulados dependen de la evolución clínica y no estarían disponibles al ingreso. \\
totcst, totmcst, income, sfdm2 & Eliminar variable & Missing/redundancia & Baja utilidad clínica directa o alta proporción de ausencia. \\
dnrday, dzclass, adlsc & Eliminar variable & Redundancia y derivación & Información duplicada o ya procesada en variables relacionadas. \\
adlp & Eliminar variable & Sesgo de recolección & Alto poder predictivo asociado al patrón de ausencia, no a una medición fisiológica estable. \\
sps, avtisst & Eliminar variable & Score derivado y leakage & Métricas procesadas o influenciadas por decisiones clínicas posteriores. \\
\end{longtable}
\endgroup

\subsubsection{Estrategia técnico-estadística (V2)}

\textbf{Filosofía y reglas técnico-estadísticas}

Esta aproximación sigue un enfoque técnico-estadístico menos dependiente del
contexto clínico. Las decisiones de preprocesamiento se basan en criterios
generales de partición, completitud de datos y medidas de tendencia central,
sin incorporar razonamiento fisiológico específico para cada variable. Por ello,
funciona como contraste metodológico frente a V1: reproduce decisiones comunes
en flujos computacionales, pero con un componente más arbitrario y menos
informado por el significado biomédico de las variables.

\textbf{Criterios de Selección y Umbrales}

\begin{itemize}
  \item Regla del 35\% de Ausencia: Bajo este enfoque, se establece un umbral
    estricto: cualquier variable con más de un 35\% de valores faltantes es
    eliminada del estudio (alb, bun, glucose, urine, adlp,
    totmcst). Se asume que el ruido introducido por la imputación masiva
    superaría el valor informativo de la variable.
  \item Eliminación por Redundancia Estadística: Se eliminan variables como dnrday
    basándose únicamente en su solapamiento matemático con otras variables, sin
    considerar su significado clínico.
  \item Prevención de Leakage Estándar: Al igual que en la V1, se eliminan las
    variables obvias de fuga de datos como death y d.time, pero se mantienen
    otras más procesadas como surv2m o surv6m bajo la premisa de que son
    predictores numéricos útiles, independientemente de su origen.
\end{itemize}

\textbf{Imputación Global}

En lugar de considerar el diagnóstico o el contexto del paciente, se aplican
técnicas de imputación masiva centralizada:

\begin{itemize}
  \item Mediana Global: Para todas las variables numéricas (wblc, bili, ph, pafi,
    charges, prg2m, etc.), se utiliza la mediana de toda la población. Este
    enfoque asume que la distribución general es suficiente para llenar los
    huecos de información, simplificando la computación del modelo.
  \item Imputación por Moda: Para las variables categóricas o de niveles (adls, edu,
    sfdm2, income), se utiliza la moda (el valor más frecuente). A diferencia de
    la V1, no se crea una categoría "Unknown", sino que se asigna al registro
    faltante la característica mayoritaria, bajo el supuesto estadístico de la
    probabilidad de mayor frecuencia.
\end{itemize}

\begingroup
\footnotesize
\setlength{\tabcolsep}{3pt}
\begin{longtable}{@{}p{0.23\textwidth}p{0.17\textwidth}p{0.21\textwidth}p{0.31\textwidth}@{}}
\caption[Preprocesamiento V2]{Preprocesamiento técnico-estadístico (V2)}\label{tab:preproc-v2}\\
\toprule
\textbf{Variable} & \textbf{Acción} & \textbf{Estrategia aplicada} & \textbf{Observación / Justificación} \\
\midrule
\endfirsthead
\multicolumn{4}{c}{\tablename\ \thetable\ — V2 (continuación)} \\
\toprule
\textbf{Variable} & \textbf{Acción} & \textbf{Estrategia aplicada} & \textbf{Observación / Justificación} \\
\midrule
\endhead
\bottomrule
\endfoot
meanbp, hrt, resp, temp, sod, scoma, sps, aps, crea, avtisst & Eliminar observaciones & Eliminación puntual & Variables casi completas; se eliminaron registros aislados con nulos. \\
race, dnr & Imputar & Categoría ``Unknown'' & Tratamiento categórico simple para ausencia administrativa. \\
charges, wblc, bili, ph, pafi, prg2m, prg6m, totcst & Imputar & Mediana global & Regla uniforme, independiente del diagnóstico. \\
adls, edu, sfdm2, income & Imputar & Moda & Sustitución por la categoría más frecuente. \\
alb, totmcst, bun, glucose, urine, adlp & Eliminar variable & Missing >35\% & Exclusión automática por umbral técnico de valores faltantes. \\
dnrday & Eliminar variable & Redundancia & Solapamiento directo con dnr. \\
surv2m, surv6m & Eliminar variable & Leakage potencial & Variables relacionadas con supervivencia futura. \\
death, d.time & Eliminar variable & Leakage directo y temporal & Información ligada al desenlace o seguimiento posterior. \\
\end{longtable}
\endgroup

\textbf{Diferencias Clave en el Tratamiento de Datos}

Mientras la estrategia clínico-estadística (V1) asume que la ausencia de datos
es informativa y dependiente del estado del paciente, la estrategia
técnico-estadística (V2) trata la ausencia como un error o vacío de información
que debe ser normalizado hacia el centro de la distribución global para no
distorsionar el aprendizaje del algoritmo.

\subsection{Preparación experimental}

Aquí vamos a desglosar el protocolo técnico y la metodología aplicados para la
construcción de los entornos experimentales. El diseño se fundamenta en la
creación de dos flujos de procesamiento diferenciados, cuya comparación
permitirá evaluar la eficiencia de la curación de datos basada en conocimiento
de dominio frente a métodos rígidos.

\subsubsection{Protocolo de partición y estratificación de la muestra}

Como etapa preliminar a cualquier intervención sobre los datos, se procedió a la
división del dataset original mediante un esquema de validación externa:

\begin{itemize}
  \item Proporción de la muestra: Se determinó una división del 80\% para el conjunto
    de entrenamiento, 10\% para el conjunto de validación y 10\% para el conjunto
    de prueba independiente.
  \item Mecanismo de estratificación: Se implementó la función createDataPartition
    sobre la variable objetivo hospdead. Esta decisión estratégica asegura que
    la prevalencia de la mortalidad hospitalaria sea equivalente en las tres
    particiones, mitigando sesgos de selección y garantizando que el modelo se
    entrene y evalúe sobre distribuciones de clase representativas del fenómeno
    clínico.
  \item Uso de los subconjuntos: El conjunto de entrenamiento se utilizó para
    ajustar los modelos y, cuando correspondía, ejecutar validación cruzada
    interna. En los experimentos con CV, la selección de hiperparámetros y del
    mejor modelo se realizó con los pliegues internos del entrenamiento, y la
    evaluación externa se efectuó sobre la unión de validación y prueba. En los
    experimentos sin CV, el conjunto de validación externo se utilizó para
    comparar configuraciones y seleccionar el mejor modelo, reservando el
    conjunto de prueba para la evaluación final.
\end{itemize}

\begin{table}[H]
  \centering
  \caption{Caracterización de los conjuntos de entrenamiento utilizados en modelado.}
  \label{tab:caracterizacion-modelado}
  \small
  \begin{tabular}{lrrrr}
    \toprule
    \textbf{Escenario} & \textbf{Observaciones} & \textbf{Predictores} & \textbf{No} & \textbf{Yes} \\
    \midrule
    V1 -- Bio Data & 7,232 & 25 & 5,372 & 1,860 \\
    V2 -- Tech Data & 7,166 & 35 & 5,311 & 1,855 \\
    V2 -- Tech Data sin columnas de riesgo & 7,166 & 26 & 5,311 & 1,855 \\
    V1 -- Bio Data con downsampling & 7,232 & 25 & 5,372 & 1,860 \\
    \bottomrule
  \end{tabular}
\end{table}

\subsubsection{Estrategia V1: procesamiento clínico-estadístico}

Este flujo se diseñó bajo la premisa de preservar la coherencia biológica y el
rigor fisiológico. Las decisiones tomadas en este entorno se detallan a
continuación:

Para la prevención de fuga de datos (Data Leakage), se ejecutó una limpieza profunda orientada a eliminar variables que introdujeran
información posterior al ingreso o directamente ligada al desenlace. Se
eliminaron columnas como death, d.time, surv2m, slos y charges. La razón fundamental es
evitar la fuga de datos (Data Leakage); la inclusión de información futura
generaría modelos con una precisión artificialmente elevada pero carentes de
utilidad en un escenario predictivo en tiempo real al ingreso del paciente.

\paragraph{Tratamiento avanzado de valores faltantes (NAs)}

Se aplicaron criterios diferenciados según la relevancia de la variable:

\begin{itemize}
  \item Criterio de exclusión por filas: Se eliminaron los registros con NAs en
    constantes vitales basales y biomarcadores críticos (meanbp, hrt, resp,
    temp, sod, scoma, aps, crea). Se asume que la ausencia de estos datos clave
    invalida la fiabilidad clínica del registro.
  \item Imputación por Mediana Agrupada: Para variables como wblc, bili y alb, se
    calculó la mediana estratificada por grupo de diagnóstico (dzgroup). Esta
    decisión avanzada responde a la lógica de que los valores basales difieren
    según la patología (ej. la bilirrubina esperada en un paciente con Cirrosis
    difiere de uno con Cáncer de Pulmón), preservando así la naturaleza clínica
    de los datos.
\end{itemize}

\paragraph{Ingeniería de características y categorización informativa}

\begin{itemize}
  \item Codificación de la educación: Se transformó la variable edu de numérica a
    categórica, definiendo niveles: "Básica" ($\leq$8 años), "Secundaria" ($\leq$12 años)
    y "Superior" (>12 años).
  \item Gestión de la incertidumbre: Variables como race y dnr con valores nulos se
    recodificaron como "Unknown". Esto permite que el algoritmo de aprendizaje
    identifique si la ausencia de registro en estas variables posee un valor
    predictivo latente.
\end{itemize}

\subsubsection{Estrategia V2: procesamiento técnico-estadístico}

Este flujo representa una aproximación técnico-estadística menos dependiente del
dominio clínico, basada en umbrales de completitud, particiones estándar y
reglas de tendencia central global.

\textbf{Criterios de limpieza y exclusión}

A diferencia de la V1, se aplicó un criterio de exclusión basado principalmente
en completitud y simplicidad operativa. Se eliminaron directamente variables con
altos volúmenes de nulidad, como alb, glucose y bun, bajo la premisa de que un
exceso de datos faltantes compromete la recuperación fiable del predictor,
optando por un modelo con menos variables pero con mayor densidad de datos
observados.

\textbf{Tratamiento de nulos mediante estrategia de imputación global}

En este escenario, el preprocesamiento se desvincula del contexto diagnóstico:

\begin{itemize}
  \item Imputación por Moda: Se desarrolló la función get\_mode para completar
    variables categóricas como edu, adls, income y sfdm2, asumiendo que el
    paciente posee el valor más frecuente de la muestra general.
  \item Imputación por Mediana Global: Para las variables numéricas restantes
    (charges, wblc, bili, ph, pafi), se utilizó la mediana de todo el conjunto
    de entrenamiento de forma centralizada, simplificando la estructura del
    modelo al eliminar la dependencia de la variable dzgroup.
  \item Consistencia de tipos: Se realizó una conversión explícita a factores para
    variables binarias leídas originalmente como numéricas (diabetes, dementia),
    asegurando que el target hospdead sea tratado como un factor con niveles "0"
    y "1".
\end{itemize}

\subsubsection{Estandarización y normalización de variables}

Para ambos entornos experimentales, la normalización se integró dentro del flujo
de modelado mediante recetas de tidymodels. Aunque durante el preprocesamiento
se conservaron versiones base y escaladas de los datos, el entrenamiento final
se realizó sobre las versiones base, delegando la estandarización al flujo de
modelado
del modelo:

\begin{itemize}
  \item Aislamiento de parámetros: Las estadísticas de normalización (media y
    desviación estándar) se calcularon exclusivamente con los datos de
    entrenamiento dentro de la receta.
  \item Prevención de contaminación: Dichos parámetros se aplicaron posteriormente
    a los conjuntos de validación y prueba. Esta decisión estratégica es
    fundamental en ciencia de datos para evitar la contaminación del modelo con
    información de los conjuntos de evaluación, garantizando la integridad de la
    comparación experimental.
  \item Codificación consistente: Las variables categóricas fueron tratadas como
    predictores nominales y transformadas a variables indicadoras dentro de la
    misma receta, incorporando además control de niveles nuevos y eliminación de
    predictores sin variabilidad.
\end{itemize}

\subsection{Marco Experimental y Selección de Algoritmos}

El diseño experimental de esta investigación se fundamenta en un enfoque de
benchmarking. El objetivo no es solo identificar el modelo con mayor
capacidad predictiva, sino evaluar cómo la calidad intrínseca del dato
(Estrategia Biológica vs. Técnica), el esquema de validación y el manejo del
desbalance impactan en dicha predicción.

\subsubsection{Selección de Modelos Supervisados}

Se han seleccionado algoritmos que representan diferentes familias de
aprendizaje automático para garantizar una comparativa exhaustiva:

\begin{itemize}
  \item Modelos Paramétricos:
    \begin{itemize}
      \item Regresión Logística (GLM): Utilizada como línea base (baseline) por su
        interpretabilidad y eficiencia en problemas de clasificación binaria.
    \end{itemize}
  \item Modelos Basados en Árboles y Reglas:
    \begin{itemize}
      \item Decision Tree (rpart): Para capturar relaciones no lineales simples y
        estructuras de decisión.
    \end{itemize}
  \item Modelos de Ensamble (Ensemble Learning):
    \begin{itemize}
      \item Random Forest (rf): Para reducir la varianza mediante bagging y mejorar
        la robustez frente a outliers.
      \item XGBoost: Para optimizar la capacidad predictiva mediante boosting
        secuencial regularizado.
    \end{itemize}
  \item Modelos Probabilísticos:
    \begin{itemize}
      \item Naive Bayes: Para incorporar una referencia probabilística simple,
        basada en independencia condicional entre predictores.
    \end{itemize}
  \item Modelos Basados en Distancia:
    \begin{itemize}
      \item K-Nearest Neighbors (KNN): Integrado dentro del mismo flujo
        normalizado, debido a su sensibilidad a la magnitud de las variables.
    \end{itemize}
\end{itemize}

\subsubsection{Estrategias de Validación y Automatización}

La experimentación se ha modularizado mediante el desarrollo de una librería de
utilidades propia (modeling\_utils\_tidymodels.R), lo que permite una ejecución sistemática y centralizada. Se definieron dos metodologías de validación:

\begin{enumerate}
  \item Validación Cruzada (5-fold Cross-Validation): Empleada para la estimación
    del error de generalización durante el ajuste inicial. Al dividir el
    conjunto de entrenamiento en 5 pliegues, se asegura que cada observación sea
    utilizada tanto para entrenamiento como para validación interna, mitigando el
    riesgo de sobreajuste (overfitting). La selección de hiperparámetros y del
    mejor modelo se realizó con las métricas internas de estos pliegues.
  \item Entrenamiento Simple (sin CV): Utilizado para evaluar el ajuste directo de
    modelos sobre el conjunto de entrenamiento y comparar sus configuraciones en
    el conjunto de validación. Esta estrategia permite observar el comportamiento
    de los modelos en su estado base y comparar la eficiencia de los modelos
    frente a la validación cruzada.
\end{enumerate}

\subsubsection{Organización computacional del proyecto}

El flujo experimental se implementó en el ecosistema R mediante archivos
RMarkdown, separando las etapas de análisis exploratorio, preprocesamiento,
entrenamiento, evaluación, test, generación metrica y comparación de en base a metricas de rendimiento. Esta organización permitió
mantener trazabilidad entre los datos originales, las transformaciones aplicadas,
los modelos entrenados, validados y testeados. La estructura detallada de
carpetas del proyecto se presenta en el Anexo~\ref{sec:anexo-estructura}.

Las principales librerías utilizadas incluyeron tidyverse, dplyr, ggplot2,
tidymodels, workflows, recipes, parsnip, tune, yardstick, rsample, pROC,
DataExplorer, corrplot, Hmisc, knitr, kableExtra, plotly, htmltools, caret y themis. La
reproducibilidad se apoyó en particiones de entrenamiento, validación y prueba
persistidas, versiones base de los datos procesados, semillas de aleatorización
y funciones centralizadas en modeling\_utils\_tidymodels.R para entrenar,
evaluar y comparar los modelos bajo criterios homogéneos.

\subsubsection{Arquitectura del Flujo de Trabajo (V1 vs. V2)}

El sistema se organiza en cuatro ejes experimentales representados por archivos
RMarkdown (.Rmd) independientes, lo que permite una trazabilidad total:

\begin{itemize}
  \item Escenario Biológico (V1 - Bio Data): Datos preprocesados bajo un criterio de
    sentido clínico, donde la imputación y selección de variables respetan la
    naturaleza médica del dataset SUPPORT2.
  \item Escenario Técnico (V2 - Tech Data): Datos procesados mediante técnicas
    automatizadas y criterios puramente estadísticos (ej. eliminación por umbral
    de valores faltantes).
\end{itemize}

En cada estrategia se ejecutaron dos sub-escenarios principales: entrenamiento
con validación cruzada y entrenamiento simple sin validación cruzada. Ambos
parten de los datos base y aplican las transformaciones necesarias dentro de la
receta de tidymodels, favoreciendo una comparación metodológicamente consistente
entre escenarios. Además, se incorporaron dos análisis complementarios: una
versión de V2 con eliminación de variables potencialmente problemáticas por
leakage o baja coherencia biológica, y un experimento V1 con downsampling dentro
de la validación cruzada para estudiar el efecto del desbalance sobre la
sensibilidad.

\subsubsection{Configuración de Hiperparámetros y Control}

Independientemente del modelo, se ha unificado el entrenamiento mediante
workflows de tidymodels, integrando receta, especificación del modelo,
búsqueda de hiperparámetros y evaluación probabilística:

\begin{itemize}
  \item Generación de probabilidades de clase: Para la construcción de curvas ROC.
  \item Optimización por ROC-AUC: Para priorizar la capacidad discriminativa global
    durante la selección de configuraciones.
  \item Registro de la mejor configuración: Para conservar los hiperparámetros del
    modelo con mayor rendimiento. En CV, este criterio se aplicó sobre las
    métricas internas de los pliegues; en entrenamiento simple, sobre el conjunto
    de validación externo.
\end{itemize}

La Tabla~\ref{tab:config-modelos} resume los hiperparámetros ajustados en cada
familia de modelos y los rangos utilizados en el entrenamiento simple sin CV.
En los escenarios con validación cruzada se mantuvo el mismo conjunto de
hiperparámetros por modelo, pero la construcción de la grilla se realizó con
\texttt{tune\_grid()} y \texttt{grid\_size = 10}, usando los rangos por defecto
definidos por tidymodels/dials para cada parámetro. En los escenarios sin CV se
empleó una grilla manual de tamaño equivalente para evaluar cada combinación
contra el conjunto de validación externo. Por esta razón, las configuraciones
óptimas obtenidas con CV pueden ubicarse fuera de los rangos manuales definidos
para el entrenamiento simple.

\begin{table}[H]
  \centering
  \caption{Configuración de hiperparámetros por modelo.}
  \label{tab:config-modelos}
  \scriptsize
  \begin{tabular}{p{0.22\textwidth}p{0.28\textwidth}p{0.40\textwidth}}
    \toprule
    \textbf{Modelo} & \textbf{Hiperparámetros ajustados} & \textbf{Grilla manual sin CV / configuración base} \\
    \midrule
    Logistic Regression & penalty, mixture & penalty: $10^{-6}$--$10^{0}$; mixture: 0--1; motor glmnet. \\
    Decision Tree & cost\_complexity, tree\_depth, min\_n & cost\_complexity: $10^{-6}$--$10^{-1}$; tree\_depth: 1--15; min\_n: 2--20; motor rpart. \\
    Random Forest & mtry, min\_n & mtry: 1--p, donde p es el número de predictores tras la receta; min\_n: 2--20; trees fijo en 500; motor ranger. \\
    Naive Bayes & smoothness, Laplace & smoothness: 0.5--1.5; Laplace: 0--3; motor klaR. \\
    XGBoost & trees, tree\_depth, learn\_rate, loss\_reduction, min\_n & trees: 200--700; tree\_depth: 2--8; learn\_rate: $10^{-4}$--$10^{-1}$; loss\_reduction: $10^{-6}$--$10^{1}$; min\_n: 2--20. \\
    KNN & neighbors, weight\_func, dist\_power & neighbors entre 3 y el máximo permitido por dimensión/tamaño muestral; weight\_func: rectangular, triangular y epanechnikov; dist\_power: 1 o 2. \\
    \bottomrule
  \end{tabular}
\end{table}

El experimento con downsampling se implementó con themis dentro de la receta de
tidymodels. Esta operación se activa únicamente en el flujo con validación
cruzada, por lo que el remuestreo se aplica a cada partición de entrenamiento
de los folds y no al dataset completo ni a los subconjuntos de evaluación.

\subsection{Métricas de Evaluación de Rendimiento}

Dado que las métricas de clasificación se describen conceptualmente en la
Sección~\ref{sec:metricas-clasificacion}, este apartado delimita únicamente su
uso dentro del protocolo experimental. Para cada modelo y escenario se
reportaron AUC-ROC, sensibilidad, especificidad, accuracy y F1-score mediante
funciones centralizadas de evaluación.

El AUC-ROC se utilizó como métrica principal para comparar la capacidad
discriminativa global entre algoritmos y estrategias de preprocesamiento. La
sensibilidad y la especificidad se emplearon para interpretar el balance entre
la detección de pacientes fallecidos y supervivientes, respectivamente. El
F1-score se utilizó como medida complementaria para valorar el equilibrio entre
precisión y recall de la clase positiva. Accuracy se mantuvo como métrica
complementaria, interpretada con cautela debido al desbalance moderado de la
variable hospdead.

Además de las métricas tabulares, se generaron curvas ROC comparativas para cada
escenario experimental. Estas curvas se presentan en las
Figuras~\ref{fig:roc-v1-cv}--\ref{fig:roc-v1-down} y permiten contrastar
visualmente el comportamiento relativo de los modelos en el espacio de
sensibilidad frente a 1-especificidad.
